\section{Avances en IoT y AAA para la gestión sostenible de residuos}

En los últimos años, la seguridad alimentaria global ha sido un tema de creciente preocupación debido a factores como el cambio climático, el crecimiento poblacional y la urbanización, los cuales han generado una presión sin precedentes sobre los sistemas alimentarios. En este contexto, la reducción y valorización de los residuos y subproductos agroalimentarios se presentan como estrategias clave para mejorar la eficiencia en el uso de los recursos y mitigar el desperdicio de alimentos. Estudios recientes destacan que la Industria 4.0 ofrece oportunidades significativas para abordar este desafío, mediante el uso de tecnologías emergentes como la inteligencia artificial, el Internet de las Cosas (IoT), blockchain, sensores inteligentes y robótica, las cuales permiten optimizar la producción, mejorar la trazabilidad y convertir los desechos de origen vegetal en productos de valor agregado \parencite{ait-kaddourTransformingPlantbasedWaste2024}.

El monitoreo de variables ambientales es fundamental para optimizar los procesos de bioconversión de residuos, especialmente en sistemas que utilizan \textit{Hermetia illucens} (BSF). Estudios han demostrado que mantener temperaturas entre $25-30\ ^{\circ}\mathrm{C}$ y niveles de humedad superiores al $60\%$ mejora significativamente la actividad larvaria y la tasa de degradación de residuos orgánicos \parencite{salamEffectDifferentEnvironmental2022}. Además, el uso de sensores IoT permite medir en tiempo real parámetros clave como gases ($CO_2$, $CH_4$) y pH, cuyos datos, al ser analizados con Algoritmos de Aprendizaje Automático (AAA), facilitan ajustes automáticos para maximizar la eficiencia del proceso y minimizar impactos ambientales. Estas tecnologías no solo contribuyen a la reducción de emisiones de gases de efecto invernadero (GEI), sino que también optimizan el uso de recursos y reducen costos operativos \parencite{liuOptimizingDataPipelines2023,vanIntegrationInternetofThingsSustainable2022}.

IoT y AAA han transformado diversos sectores, como la agricultura, donde las redes de sensores inalámbricos han optimizado el riego y el uso de fertilizantes \parencite{satputeREVIEWPAPERWIRELESS2021}, la manufactura, donde el análisis predictivo ha reducido tiempos de inactividad y mejorado la calidad de los productos \parencite{avalekarOptimizingAgriculturalEfficiency2024}, y la salud, donde dispositivos \textit{wearables} han revolucionado el diagnóstico y seguimiento de pacientes \parencite{paulMachineLearningIoT2024}. Dado el éxito de IoT y AAA en estos sectores, su aplicación en la bioconversión de residuos mediante \textit{Hermetia illucens} representa una oportunidad significativa para mejorar la sostenibilidad y la eficiencia en la gestión de residuos agroindustriales.

Diversos estudios han demostrado que el monitoreo en tiempo real de parámetros como temperatura, humedad y composición de gases permite optimizar el crecimiento de las larvas y reducir la emisión de GEI durante el proceso de bioconversión \parencite{rossiEstimatingDynamicsGreenhouse2024}. En particular, la estimación de las dinámicas de emisión de CO$_2$ y N$_2$O bajo condiciones controladas ha evidenciado cómo la regulación ambiental mediante sensores inteligentes puede minimizar las pérdidas de carbono y mejorar la eficiencia del proceso. 

Un análisis de la huella de carbono y agua en la producción de \textit{Hermetia illucens} en una granja en España mostró que la implementación de estrategias de monitoreo y control automatizado reduce significativamente el impacto ambiental, favoreciendo una producción más sostenible y alineada con los principios de la economía circular \parencite{IPCC2014,galan-diazCarbonWaterFootprint2024}. Además, el tratamiento de residuos orgánicos como el estiércol avícola con tecnología basada en \textit{Hermetia illucens} contribuye a la reducción de emisiones de N$_2$O y CO$_2$ en los suelos, lo que refuerza su potencial como estrategia clave para la mitigación del cambio climático en sistemas agrícolas y pecuarios \parencite{jenkinsProcessingPoultryManure2025}.

Aunque su implementación aún enfrenta desafíos tecnológicos y económicos, la integración de IoT y AAA en estos sistemas representa una oportunidad única para desarrollar soluciones innovadoras que optimicen el uso de recursos y promuevan la neutralidad de carbono en la producción agroalimentaria.